% simple-minimizer-proof.tex — \input'd from main.tex
% Conventions: thesis/CLAUDE.md (audiences, content rules, format rules, comments)
%
% Sources: Jörn's 120-minute talk (January 2025) expanding the 1-page
%   sketch in Haim-Kislev 2017.
%
% Structure:
%   5.1 Primal-Dual Equivalence (Clarke's Dual Action Principle for polytopes)
%   5.2 Simplification Algorithm (5 steps: piecewise linear, pure Reeb,
%       coalesce, rescale, compactness)

% No \section here: this file is \input'd inside Section "Simple Minimizer
% Existence" in chapter-algorithm.tex, so its subsections become 5.x.

The proof follows \cite{HK2017} but with complete details. We first
establish a primal-dual equivalence (Clarke's Dual Action Principle
adapted to polytopes), then construct a simple minimizer via a
five-step simplification algorithm.

%%% =================================================================
%%% Subsection 5.1: Primal-Dual Equivalence
%%% =================================================================

\subsection{Primal-Dual Equivalence}\label{subsec:primal-dual}

% QC: Clarke's Dual Action Principle is classical for smooth convex bodies
%   (Clarke 1979, 1981). The polytope case requires care with the gauge
%   function (not smooth) and the dual functional (integral of h_K^2).
%   Literature on this is sparse and often omits details. This writeup
%   is self-contained for a master's thesis audience.
% Downstream: establishes that minimizing action over Reeb orbits is
%   equivalent to minimizing I_K over centered curves. The dual problem
%   is easier to work with for the simplification algorithm.

We work with absolutely continuous curves
$z \in W^{1,2}([0, T], \mathbb{R}^4)$
(the Sobolev space $H^1$, which embeds into $C^0$).
% QC: W^{1,2} functions on a compact interval are absolutely continuous
%   and continuous. No need for a separate C^0 assumption.

Recall the gauge function $g_K(x) = \inf\{\lambda > 0 : x \in \lambda K\}$
(Definition~\ref{def:gauge-function}) and the support function
$h_K(v) = \max_{x \in K} \langle x, v \rangle$
(Definition~\ref{def:support-function}).
The Hamiltonian is $H = g_K^2$ (Definition~\ref{def:hamiltonian}).

\begin{lemma}[Fenchel duality]\label{lem:fenchel-duality}
% QC: Legendre-Fenchel conjugate of g_K^2 is (1/4)h_K^2.
%   The pointwise inequality is Young's inequality for this pair.
%   Equality characterizes the subgradient relation.
% Downstream: integrated along curves, this links the primal
%   (action + gauge constraint) to the dual (I_K functional).
The functions $g_K^2$ and $\tfrac{1}{4}h_K^2$ are Legendre-Fenchel conjugates:
\begin{equation}\label{eq:conjugacy}
  g_K^2(x) = \sup_y \bigl\{\langle x, y \rangle - \tfrac{1}{4} h_K^2(y)\bigr\},
  \qquad
  \tfrac{1}{4}h_K^2(y) = \sup_x \bigl\{\langle x, y \rangle - g_K^2(x)\bigr\}.
\end{equation}
Consequently, for any $x, y \in \mathbb{R}^4$:
\begin{equation}\label{eq:fenchel-inequality}
  g_K^2(x) + \tfrac{1}{4} h_K^2(y) \geq \langle x, y \rangle,
\end{equation}
with equality if and only if
$y \in \partial g_K^2(x)$ if and only if
$x \in \partial(\tfrac{1}{4} h_K^2)(y)$.
\end{lemma}

\begin{proof}
% Jörn: good enough though not polished
\emph{Conjugacy.}
We compute the Legendre-Fenchel conjugate of~$g_K^2$ directly.
Parametrize $x = \lambda u$ with $\lambda \geq 0$ and $g_K(u) = 1$,
so $g_K^2(x) = \lambda^2$.
Then:
\[
  (g_K^2)^*(y)
  = \sup_x \bigl\{\langle x, y \rangle - g_K^2(x)\bigr\}
  = \sup_{\lambda \geq 0}\;
    \sup_{g_K(u) = 1}
    \bigl\{\lambda \langle u, y \rangle - \lambda^2\bigr\}
  = \sup_{\lambda \geq 0}
    \bigl\{\lambda\, h_K(y) - \lambda^2\bigr\},
\]
where we used $\sup_{g_K(u) = 1} \langle u, y \rangle = h_K(y)$
(the supremum of a linear functional over a convex body
is attained on its boundary).
The quadratic $\lambda h_K(y) - \lambda^2$ is maximized
at $\lambda = h_K(y)/2$, giving
\[
  (g_K^2)^*(y) = \tfrac{1}{4} h_K^2(y).
\]
This is the second relation in~\eqref{eq:conjugacy}.
The first follows by the same computation with $K$ replaced by its
polar body $K^\circ$ (whose gauge function is $h_K$ and
support function is $g_K$).

\emph{Inequality.}
By definition of the conjugate,
$(g_K^2)^*(y) \geq \langle x, y \rangle - g_K^2(x)$ for any~$x$.
Substituting $(g_K^2)^* = \tfrac{1}{4}h_K^2$ and rearranging
gives~\eqref{eq:fenchel-inequality}.

\emph{Equality condition.}
Equality holds iff $x$ attains the supremum in $(g_K^2)^*(y) = \sup_z \{\langle z, y \rangle - g_K^2(z)\}$,
i.e., $x$ maximizes $z \mapsto \langle z, y \rangle - g_K^2(z)$.
By the first-order optimality condition, this is equivalent to
$y \in \partial g_K^2(x)$.
By the same argument applied to the first relation in~\eqref{eq:conjugacy},
it is equivalent to $x \in \partial(\tfrac{1}{4}h_K^2)(y)$.
\end{proof}

\begin{definition}[Primal problem]\label{def:primal-problem}
% Jörn: text approved (05aa8b3)
% QC: restatement of the EHZ capacity minimization, using H = g_K^2.
%   The Hamiltonian formulation (differential inclusion with subdifferential)
%   is equivalent to the Reeb orbit formulation on polytopes.
The \emph{primal functional} is the action
\[
  A(\gamma) = \tfrac{1}{2}\int_0^T
  \langle -J_0\dot{\gamma},\, \gamma \rangle\, dt.
\]
The \emph{primal problem} is:
\begin{multline}\label{eq:primal-problem}
  c_{\mathrm{EHZ}}(K) = \min \Bigl\{ A(\gamma) :
  \gamma \in W^{1,2}([0,T], \mathbb{R}^4),\;
  \gamma(0) = \gamma(T),\;
  T > 0, \\
  g_K^2(\gamma) \equiv 1,\;
  -J_0\dot{\gamma} \in \partial g_K^2(\gamma) \text{ a.e.} \Bigr\}.
\end{multline}
The constraints say: $\gamma$ lies on $\partial K$
(the level set $\{g_K^2 = 1\}$) and its velocity satisfies the
Hamiltonian differential inclusion, which equals the Reeb orbit
differential inclusion (Definition~\ref{def:simple-reeb-orbit}).
The action equals the period: $A(\gamma) = T$.
\end{definition}

We now introduce an equivalent dual problem without motivation.

\begin{definition}[Dual functional and dual problem]\label{def:dual-problem}
% Jörn: text approved (05aa8b3)
% QC: I_K depends only on the velocity dz/dt, not on z itself.
%   This is the key advantage of the dual formulation: no position constraint
%   (gamma on boundary K) and no differential inclusion constraint.
%   The centering condition removes the translation ambiguity in going
%   from gamma to z = gamma - center(gamma).
The \emph{dual functional} is
\[
  I_K(z) = \frac{1}{4} \int_0^T h_K^2\bigl( -J_0 \dot{z}(t) \bigr) \, dt.
\]
The \emph{dual problem} is:
\begin{multline}\label{eq:dual-problem}
  \min \Bigl\{ I_K(z) :
  z \in W^{1,2}([0,T], \mathbb{R}^4),\;
  z(0) = z(T),\;
  T > 0, \\
  A(z) = T,\;
  \textstyle\frac{1}{T} \int_0^T z = 0 \Bigr\}.
\end{multline}
The constraint $\frac{1}{T} \int z = 0$ is the
\emph{centering condition}; it removes the translation
ambiguity (since $I_K$ depends only on $\dot{z}$, not on~$z$).
\end{definition}

\cite{HK2017} also states the dual problem without motivation
(in a slightly different form).
The dual functional and constraints can be motivated by working
the equivalence proof
(Theorem~\ref{thm:primal-dual-equivalence}) in reverse.

\begin{theorem}[Primal-dual equivalence (Clarke's Dual Action Principle)]%
\label{thm:primal-dual-equivalence}
% QC: the critical-point correspondence is the core content.
%   Existence of minimizers is a separate fact (thm:orbit-existence-smooth/polytope
%   for the primal; compactness + l.s.c. for the dual).
%   The equivalence proof does not use existence of minimizers.
% Downstream: allows us to work in the dual formulation, where the
%   simplification algorithm (subsec:simplification) is natural.
Let $K$ be a convex body. Critical points of the primal and
dual problems correspond bijectively via translation:
if $\gamma$ is a primal critical point (Hamiltonian orbit on
$\partial K$), then
$z = \gamma - \frac{1}{T} \int_0^T \gamma(t) \, dt$
is a dual critical point, and conversely.
At every corresponding pair, the values are equal:
\begin{equation}\label{eq:primal-dual-values}
  I_K(z) = T = A(\gamma).
\end{equation}
Consequently, $\min_\gamma A(\gamma) = \min_z I_K(z) = c_{\mathrm{EHZ}}(K)$.
\end{theorem}

\begin{proof}
% Jörn: good enough though not polished
\emph{Part 1: Euler--Lagrange characterization of dual critical points.}
We characterize critical points of the dual functional
$I_K(z) = \frac{1}{4}\int_0^T h_K^2(-J_0\dot{z}) \, dt$
under the constraints:
\begin{itemize}
\item closure: $z(0) = z(T)$,
\item action: $A(z) = T$,
\item centering: $\int_0^T z(t) \, dt = 0$.
\end{itemize}

We use the calculus of variations with Lagrange multipliers.
Form the Lagrangian incorporating the constraints:
\[
  L(z, \dot{z}) = \tfrac{1}{4}h_K^2(-J_0\dot{z}) + \alpha\langle -J_0\dot{z}, z\rangle + \langle \mu, z\rangle,
\]
where $\alpha \in \mathbb{R}$ is the multiplier for the action constraint
and $\mu \in \mathbb{R}^4$ for the centering constraint.

The Lagrangian is convex (but not differentiable) in $\dot{z}$
and smooth (linear) in~$z$.
For such Lagrangians, a critical point $z$ satisfies the
\emph{Euler--Lagrange inclusion}: there exists an absolutely continuous
$\xi(t) \in \partial_{\dot{z}} L(z(t), \dot{z}(t))$ a.e.\ such that
\begin{equation}\label{eq:euler-lagrange-inclusion}
  \dot{\xi}(t) = \nabla_z L(z(t), \dot{z}(t)) \quad\text{a.e.}
\end{equation}
(see Clarke, \emph{Optimization and Nonsmooth Analysis}, Ch.~5).

\emph{Right-hand side of~\eqref{eq:euler-lagrange-inclusion}.}
Since $\tfrac{1}{4}h_K^2(-J_0\dot{z})$ is independent of~$z$,
only the multiplier terms contribute:
\[
  \nabla_z L = -\alpha J_0 \dot{z} + \mu.
\]

\emph{Left-hand side of~\eqref{eq:euler-lagrange-inclusion}.}
The subdifferential of $L$ in $\dot{z}$ is
% chain rule: ∂_ż[f(Aż)] = Aᵀ ∂f(Aż) with A = -J₀, so Aᵀ = J₀
\[
  \partial_{\dot{z}} L
  = J_0\, \partial\bigl(\tfrac{1}{4}h_K^2\bigr)(-J_0\dot{z})
    + \alpha J_0 z,
\]
using $(-J_0)^\top = J_0$ in the chain rule for $f \circ (-J_0)$.
So the E-L inclusion says: there exists $p(t) \in \partial(\tfrac{1}{4}h_K^2)(-J_0\dot{z}(t))$
a.e.\ such that $\xi(t) = J_0 p(t) + \alpha J_0 z(t)$ is absolutely continuous and
\[
  J_0 \dot{p} + \alpha J_0 \dot{z}
  = -\alpha J_0 \dot{z} + \mu,
  \qquad\text{i.e.,}\qquad
  J_0 \dot{p} = -2\alpha J_0 \dot{z} + \mu.
\]

Integrating over $[0,T]$ and using periodicity of~$p$ and~$z$:
\[
  J_0\underbrace{[p(T) - p(0)]}_{= 0}
  = -2\alpha J_0 \underbrace{[z(T) - z(0)]}_{= 0} + T\mu,
\]
hence $\mu = 0$. The equation $\dot{p} = -2\alpha \dot{z}$ a.e.\ integrates to
\[
  p(t) = c - 2\alpha\, z(t)
\]
for a constant $c \in \mathbb{R}^4$.
Since $p(t) \in \partial(\tfrac{1}{4}h_K^2)(-J_0\dot{z}(t))$:
\[
  c - 2\alpha\, z(t)
  \in \partial\bigl(\tfrac{1}{4}h_K^2\bigr)\bigl(-J_0\dot{z}(t)\bigr)
  \quad \text{a.e.}
\]
Writing $\nu = -2\alpha$ and $b = c$:
\begin{equation}\label{eq:dual-euler-lagrange}
  \nu z(t) + b \in
  \partial\bigl(\tfrac{1}{4} h_K^2\bigr)\bigl(-J_0\dot{z}(t)\bigr)
  \quad \text{a.e.}
\end{equation}
for some $\nu \in \mathbb{R}$ and constant $b \in \mathbb{R}^4$.

\medskip
\emph{Part 2: Primal critical point $\Rightarrow$ dual critical point.}
Let $\gamma$ be a primal critical point (Hamiltonian orbit on $\partial K$),
so $g_K^2(\gamma) \equiv 1$ and $-J_0\dot{\gamma} \in \partial g_K^2(\gamma)$ a.e.
Define
\[
  z(t) := \gamma(t) - \tfrac{1}{T}\!\int_0^T \gamma(s)\, ds.
\]
We show $z$ is a dual critical point.

\emph{Closure.}
Since $\gamma(0) = \gamma(T)$, we have $z(0) = z(T)$.

\emph{Centering.}
By definition, $\int_0^T z = 0$.

\emph{Action.}
Since $\dot{z} = \dot{\gamma}$ and $\int_0^T z = 0$:
\[
  A(z) = \tfrac{1}{2}\!\int_0^T \langle -J_0\dot{z},\, z \rangle\, dt
       = \tfrac{1}{2}\!\int_0^T \langle -J_0\dot{\gamma},\, \gamma \rangle\, dt
       = A(\gamma) = T.
\]

\emph{Euler--Lagrange equation.}
The primal differential inclusion $-J_0\dot{\gamma} \in \partial g_K^2(\gamma)$ a.e.\
is the equality case of Fenchel duality (Lemma~\ref{lem:fenchel-duality}),
which is equivalent to
\[
  \gamma \in \partial\bigl(\tfrac{1}{4}h_K^2\bigr)(-J_0\dot{\gamma})
  \quad\text{a.e.}
\]
Since $\dot{z} = \dot{\gamma}$, this gives
\[
  \gamma \in \partial\bigl(\tfrac{1}{4}h_K^2\bigr)(-J_0\dot{z})
  \quad\text{a.e.}
\]
Write $\gamma = z + b$ where $b = \frac{1}{T}\int_0^T \gamma$ is constant.
This is exactly the E-L condition~\eqref{eq:dual-euler-lagrange} with $\nu = 1$.

\medskip
\emph{Part 3: Dual critical point $\Rightarrow$ primal critical point.}
Let $z$ be a dual critical point, so $z$ satisfies the E-L
condition~\eqref{eq:dual-euler-lagrange} for some $\nu \in \mathbb{R}$ and $b \in \mathbb{R}^4$.
Define $\gamma(t) := \nu z(t) + b$.

\emph{Step 1: Converting the E-L condition.}
Since $\dot{\gamma} = \nu\dot{z}$, the E-L condition becomes
\[
  \gamma \in \partial\bigl(\tfrac{1}{4}h_K^2\bigr)(-J_0\dot{\gamma}/\nu)
  \quad\text{a.e.}
\]
By Fenchel duality (Lemma~\ref{lem:fenchel-duality}), this is equivalent to
\[
  -J_0\dot{\gamma}/\nu \in \partial g_K^2(\gamma)
  \quad\text{a.e.}
\]

\emph{Step 2: Proving $\nu = 1$.}
We show $\nu = 1$ by deriving three identities and combining them.

\emph{(a) Constancy of $g_K^2(\gamma)$.}
Since $-J_0\dot{\gamma}/\nu \in \partial g_K^2(\gamma)$ a.e.,
the subgradient inequality gives
\[
  \frac{d}{dt} g_K^2(\gamma)
  \geq \bigl\langle \tfrac{-J_0\dot{\gamma}}{\nu},\, \dot{\gamma} \bigr\rangle
  = -\frac{1}{\nu}\, \omega_0(\dot{\gamma}, \dot{\gamma}) = 0
  \quad\text{a.e.}
\]
So $t \mapsto g_K^2(\gamma(t))$ is non-decreasing.
Since $\gamma$ is periodic, a non-decreasing periodic function is constant:
$g_K^2(\gamma) \equiv c^2$ for some constant $c \geq 0$.

\emph{(b) Action of $\gamma$ via homogeneity.}
For $y \in \partial g_K^2(x)$, the subgradient inequality at $z = \lambda x$ gives
$\lambda^2 g_K^2(x) \geq g_K^2(x) + (\lambda - 1)\langle y, x\rangle$.
Dividing by $(\lambda - 1)$ and taking $\lambda \to 1$ from both sides
yields $\langle y, x \rangle = 2\, g_K^2(x)$.
Applied with $y = -J_0\dot{\gamma}/\nu$ and $x = \gamma$:
\begin{equation}\label{eq:euler-identity-action}
  \frac{1}{\nu}\langle -J_0\dot{\gamma},\, \gamma \rangle = 2c^2,
  \quad\text{so}\quad
  A(\gamma) = \nu c^2 T.
\end{equation}

\emph{(c) Action of $\gamma$ computed directly.}
Since $\gamma = \nu z + b$, $\dot{\gamma} = \nu \dot{z}$, and $\int_0^T z = 0$:
\begin{align*}
  A(\gamma)
  &= \tfrac{1}{2}\!\int_0^T \langle -J_0 \dot{\gamma},\, \gamma \rangle\, dt
   = \tfrac{1}{2}\!\int_0^T
    \bigl\langle -\nu J_0 \dot{z},\, \nu z + b \bigr\rangle\, dt \\
  &= \nu^2 A(z) + \tfrac{\nu}{2}
    \bigl\langle b,\, \underbrace{-J_0[z(T) - z(0)]}_{= 0} \bigr\rangle
   = \nu^2 A(z).
\end{align*}
Since $z$ is a dual critical point with $A(z) = T$, we have $A(\gamma) = \nu^2 T$.

\emph{(d) Integrated Fenchel equality.}
The Fenchel equality (Lemma~\ref{lem:fenchel-duality}) at
$-J_0\dot{\gamma}/\nu \in \partial g_K^2(\gamma)$ gives pointwise:
\[
  g_K^2(\gamma) + \tfrac{1}{4}h_K^2(-J_0\dot{\gamma}/\nu)
  = \langle \gamma,\, -J_0\dot{\gamma}/\nu \rangle.
\]
Since $\dot{\gamma} = \nu\dot{z}$, integrating over $[0,T]$ gives
\[
  c^2 T + I_K(z) = 2\nu T,
\]
hence $c^2 = 2\nu - 1$.

\emph{Combining (a)--(d).}
From (b) and (c): $\nu c^2 T = \nu^2 T$, i.e., $\nu(\nu - c^2) = 0$.
From (d): $c^2 = 2\nu - 1$.
If $\nu = 0$, then $c^2 = -1$, which is impossible.
So $c^2 = \nu$, and substituting: $\nu = 2\nu - 1$, hence $\nu = 1$.

\emph{Step 3: Conclusion.}
With $\nu = 1$, we have $c = 1$, so $g_K^2(\gamma) \equiv 1$
(i.e., $\gamma$ lies on $\partial K$) and
\[
  -J_0\dot{\gamma} \in \partial g_K^2(\gamma) \quad\text{a.e.}
\]
(i.e., $\gamma$ is a Hamiltonian orbit on $\partial K$).
Since $\gamma = z + b$ and $\int_0^T z = 0$,
we have $b = \mathrm{center}(\gamma)$, i.e.,
$z = \gamma - \mathrm{center}(\gamma)$.

\medskip
\emph{Part 4: Value equality.}
Let $\gamma$ be a primal critical point and $z = \gamma - \mathrm{center}(\gamma)$
the corresponding dual critical point.
The Fenchel equality (Lemma~\ref{lem:fenchel-duality}) at the subgradient
relation $-J_0\dot{\gamma} \in \partial g_K^2(\gamma)$ gives pointwise:
\[
  g_K^2(\gamma) + \tfrac{1}{4} h_K^2(-J_0\dot{\gamma})
  = \langle \gamma, -J_0\dot{\gamma} \rangle.
\]
Integrating over $[0,T]$ and using $g_K^2(\gamma) \equiv 1$:
\[
  T + I_K(\gamma) = 2A(\gamma).
\]
Since $A(\gamma) = T$ (the action equals the period for any Hamiltonian orbit),
we have $I_K(\gamma) = T$. Since $\dot{z} = \dot{\gamma}$, we get $I_K(z) = I_K(\gamma) = T$.
Therefore $I_K(z) = T = A(\gamma)$.

The correspondence is bijective, so minimizers correspond to minimizers:
$\min_\gamma A(\gamma) = \min_z I_K(z) = c_{\mathrm{EHZ}}(K)$.
\end{proof}

%%% =================================================================
%%% Subsection 5.2: Simplification Algorithm
%%% =================================================================

\subsection{Simplification Algorithm}\label{subsec:simplification}

% QC: the heart of the proof. Takes an arbitrary dual minimizer z_0
%   and constructs a simple dual minimizer z_* via 5 steps. Each step
%   is designed to preserve or improve the optimality while adding
%   structure. The final compactness argument (step 5) ensures we
%   can take a limit that is both simple and optimal.

% AI agents: this is a 1:1 restatement of thm:simple-minimizer for the reader's convenience.
\begin{theorem*}[Theorem~\ref{thm:simple-minimizer} restated]
Let $K$ be a polytope.
Among all generalized Reeb orbits on $\partial K$, there exists one
with minimum action that is simple
(Definition~\ref{def:simple-reeb-orbit}).
Consequently,
\[
  c_{\mathrm{EHZ}}(K) = \min\bigl\{ A(\gamma) :
  \gamma \text{ is a simple Reeb orbit on } \partial K \bigr\}.
\]
\end{theorem*}

\begin{proof}[Proof of Theorem~\ref{thm:simple-minimizer}]
A minimum-action generalized Reeb orbit on~$\partial K$ exists
(Theorem~\ref{thm:orbit-existence-polytope}): the minimum of the action over
all generalized Reeb orbits is attained and positive.
By the primal-dual equivalence (Theorem~\ref{thm:primal-dual-equivalence}),
any such primal minimizer~$\gamma$ corresponds to a dual minimizer
$z_0 = \gamma - \mathrm{center}(\gamma)$ with $I_K(z_0) = T = A(\gamma)$.
We will construct a simple dual minimizer~$z_*$, which then gives a
simple primal minimizer via $\gamma_* = z_* + b$.

Let $z_0$ be an arbitrary dual minimizer with period~$T > 0$. By the
primal-dual equivalence, $A(z_0) = T = I_K(z_0)$. We construct a sequence
of curves $z_0 \to z_1 \to z_2 \to z_3 \to z_4 \to z_*$, where each
arrow represents one simplification step. The final curve~$z_*$ is
simple (piecewise linear with pure Reeb vectors as velocities, each
Reeb vector used at most once) and satisfies $A(z_*) = T = I_K(z_*)$,
hence is also a minimizer.

\paragraph{Step 0: Setup.}
% QC: z_0 is any dual minimizer. We know A(z_0) = T = I_K(z_0) from the
%   primal-dual equivalence. The goal is to perturb z_0 into a simple
%   curve while preserving optimality.
We have $z_0 \in W^{1,2} \cap C^0$ with $z_0(0) = z_0(T)$,
$\frac{1}{T} \int_0^T z_0 = 0$, and $A(z_0) = T = I_K(z_0)$.

\paragraph{Step 1: Piecewise linear approximation.}
By the primal-dual correspondence (Theorem~\ref{thm:primal-dual-equivalence}),
$z_0$ corresponds to a primal minimizer $\gamma$ on $\partial K$.
Since $\gamma$ is a generalized Reeb orbit on $\partial K$, we have
$\dot{\gamma}(t) \in J_0 N_K(\gamma(t))$ a.e., which for polytopes equals
$\mathrm{conv}\{R_1,\ldots,R_F\}$ (the convex hull of facet Reeb vectors).
By the correspondence $z_0 = \gamma - \mathrm{center}(\gamma)$, we have
$\dot{z}_0(t) \in \mathrm{conv}\{R_1,\ldots,R_F\}$ a.e.

% [AGENT-WRITTEN, UNREVIEWED BY JÖRN]
% QC: translates HK2017 Lemma 4.2 (lines 275-293). The velocity-in-convex-hull
%   argument uses the integral mean value in a convex set, which is standard.
\textbf{Claim (Piecewise affine approximation):}
For any $\varepsilon > 0$, there exists a piecewise affine curve $z_1 : [0,T] \to \mathbb{R}^4$ such that:
\begin{itemize}
\item $z_1(0) = z_1(T)$ (closed),
\item $\|z - z_1\|_{W^{1,2}} < \varepsilon$ (approximation in Sobolev norm),
\item $\dot{z}'(t) \in \mathrm{conv}\{R_1,\ldots,R_F\}$ on each affine piece,
\item $|A(z_1) - A(z)| < \varepsilon$ and $|I_K(z_1) - I_K(z)| < \varepsilon$
      (by continuity of $A$ and $I_K$ in $W^{1,2}$).
\end{itemize}
The centering condition $\frac{1}{T}\int_0^T z_1 = 0$ is ensured by translating $z_1$.

\emph{Proof of Claim.}
Choose a partition $0 = t_0 < t_1 < \cdots < t_m = T$
fine enough that the piecewise affine interpolant $z_1$
defined by $z_1(t_i) = z(t_i)$ (affine on each $(t_i, t_{i+1})$)
satisfies $\|z - z_1\|_{W^{1,2}} < \varepsilon$.
(This is a standard density result for $W^{1,2}$.)
On each subinterval, the velocity is the integral mean of~$\dot{z}$:
\[
  \dot{z}'(t)
  = \frac{z(t_{i+1}) - z(t_i)}{t_{i+1} - t_i}
  = \frac{1}{t_{i+1} - t_i}\int_{t_i}^{t_{i+1}} \dot{z}(s)\, ds.
\]
Since $\dot{z}(s) \in \mathrm{conv}\{R_1,\ldots,R_F\}$ a.e.\
and the convex hull is closed and convex, the integral mean
also lies in $\mathrm{conv}\{R_1,\ldots,R_F\}$.
Closure and action/functional estimates follow from
$z_1(0) = z(0) = z(T) = z_1(T)$ and
the continuity of $A$ and $I_K$ in the $W^{1,2}$ norm.
\hfill$\diamond$

\paragraph{Step 2: Split into pure Reeb vectors.}
After Step~1, each velocity segment $\dot{z}'(t)$ is a convex combination
of facet Reeb vectors. We replace each such segment by a concatenation
of segments with \emph{pure} facet Reeb vectors (not convex combinations).

% [AGENT-WRITTEN, UNREVIEWED BY JÖRN]
% QC: action identity translates HK2017 Proposition 4.3 (lines 301-315).
%   Splitting translates HK2017 Lemma 4.4 (lines 317-354).
%   The sign-flip argument for choosing ordering is the core of Lemma 4.4's proof.
We first establish an action identity for piecewise constant velocities.

\textbf{Claim (Action identity):}
Let $z$ be a closed piecewise affine curve with
$\dot{z}(t) = \sum_{i=1}^m \mathbf{1}_{I_i}(t)\, w_i$,
where $(I_i)_{i=1}^m$ partitions $[0,T]$. Then
\begin{equation}\label{eq:action-identity}
  2A(z)
  = \int_0^T \langle -J_0 \dot{z},\, z \rangle\, dt
  = \sum_{i=1}^m \sum_{j=1}^{i-1}
    |I_i|\, |I_j|\, \omega_0(w_i, w_j).
\end{equation}

\emph{Proof of Claim.}
Write $z(t) = z(0) + \int_0^t \dot{z}(s)\, ds$.
The term $\langle -J_0 \dot{z}(t),\, z(0) \rangle$
integrates to $\langle -J_0[z(T) - z(0)],\, z(0) \rangle = 0$
by closure. So
\[
  \int_0^T \langle -J_0 \dot{z}(t),\, z(t) \rangle\, dt
  = \int_0^T \Bigl\langle -J_0 \dot{z}(t),\,
    \int_0^t \dot{z}(s)\, ds \Bigr\rangle dt.
\]
For $t \in I_i = (\tau_{i-1}, \tau_i)$, we have $\dot{z}(t) = w_i$ and
$\int_0^t \dot{z}(s)\, ds
  = \sum_{j=1}^{i-1} |I_j|\, w_j + (t - \tau_{i-1})\, w_i$.
The $w_i$ self-pairing vanishes:
$\omega_0(w_i, w_i) = 0$.
Thus
\[
  \int_{I_i} \langle -J_0 w_i,\, \textstyle\sum_{j<i} |I_j|\, w_j \rangle\, dt
  = |I_i| \sum_{j=1}^{i-1} |I_j|\, \omega_0(w_i, w_j).
\]
Summing over $i$ gives~\eqref{eq:action-identity}.
\hfill$\diamond$

\medskip

\textbf{Claim (Splitting convex combinations):}
If $\dot{z}'(t) = \sum_{i=1}^k a_i R_i$ with $a_i > 0$, $\sum a_i = 1$
on an interval~$I$,
replace $I$ by $k$ subintervals of lengths $\Delta t_i = a_i |I|$,
each with pure velocity $R_i$.
Total time preserved: $\sum \Delta t_i = |I|$.
The ordering can be chosen so that $A(z_2) \geq A(z_1)$.

\emph{Proof of Claim.}
The replacement preserves closure
($\sum \Delta t_i\, R_i = |I| \sum a_i R_i = |I|\, \dot{z}'|_I$,
same displacement).
By the action identity~\eqref{eq:action-identity},
the action change from splitting interval~$I$ depends on the
internal ordering of the pure segments. Reversing the ordering
of the $k$ pure segments within~$I$ flips the sign of
$\sum_{r < s} a_r a_s |I|^2\, \omega_0(R_s, R_r)$
(since $\omega_0(R_s, R_r) = -\omega_0(R_r, R_s)$).
Hence one of the two orderings yields
$\Delta A \geq 0$. Apply this to each affine piece.
\hfill$\diamond$

After splitting, the dual functional equals the period \emph{exactly}:

The identity $I_K(z_2) = T$ follows from the Reeb normalization:
for $\dot{z}(t) = R_i = \frac{2}{h_i} J_0 n_i$, we have
\begin{align*}
  h_K^2(-J_0 \dot{z}(t))
  &= h_K^2\Bigl( -J_0 \cdot \frac{2}{h_i} J_0 n_i \Bigr)
  = h_K^2\Bigl( \frac{2}{h_i} n_i \Bigr) \\
  &= \Bigl( \frac{2}{h_i} \Bigr)^2 h_K^2(n_i)
  = \frac{4}{h_i^2} \cdot h_i^2
  = 4,
\end{align*}
hence $I_K(z_2) = \frac{1}{4} \int_0^T 4 \, dt = T$.

\paragraph{Step 3: Merge repeated Reeb vectors.}
After Step~2, the velocity $\dot{z}''$ takes values in the finite set
$\{R_1,\ldots,R_F\}$ (pure facet Reeb vectors).
If the same Reeb vector $R_i$ appears in two disjoint time intervals,
we merge them into a single contiguous segment.

% [AGENT-WRITTEN, UNREVIEWED BY JÖRN]
% QC: translates HK2017 Lemma 4.5 (lines 356-397). The action difference
%   formula follows from the action identity by tracking which terms change.
\textbf{Claim (Merging repeated velocities):}
Suppose $w_r = w_s = R_A$ for some $r < s$.
We can merge the two $R_A$ segments into one contiguous segment
with $A(z_3) \geq A(z_2)$.

\emph{Proof of Claim.}
\textbf{Forward merge:} erase $I_s$ and extend $I_r$ by $|I_s|$,
shifting $I_{r+1}, \ldots, I_{s-1}$ later.
By the action identity~\eqref{eq:action-identity},
the only terms that change involve pairs $(i,j)$ with
$r < i < s$ and $j \in \{r, s\}$:
the term $|I_s|\,|I_i|\,\omega_0(w_s, w_i)$ (with $i < s$)
is replaced by $|I_s|\,|I_i|\,\omega_0(w_i, w_r)$ (with $i > r$).
Since $w_r = w_s$, the action difference is
\[
  \Delta A_{\text{fwd}}
  = \sum_{i=r+1}^{s-1} |I_s|\, |I_i|\, \omega_0(w_i, w_s).
\]
\textbf{Backward merge:} erase $I_r$ and extend $I_s$ by $|I_r|$.
The same calculation gives
$\Delta A_{\text{bwd}}
  = \sum_{i=r+1}^{s-1} |I_r|\, |I_i|\, \omega_0(w_r, w_i)
  = -\sum_{i=r+1}^{s-1} |I_r|\, |I_i|\, \omega_0(w_i, w_r)$.
Since $w_r = w_s$, the two differences have opposite signs
(up to the factor $|I_s|$ vs.\ $|I_r|$; when $|I_r| = |I_s|$
they are exactly opposite). In either case,
$\max(\Delta A_{\text{fwd}}, \Delta A_{\text{bwd}}) \geq 0$.
Choose that option.
Closure is preserved (total displacement unchanged).
Iterate until each $R_i$ appears at most once.
\hfill$\diamond$

The dual functional $I_K$ is unchanged: it depends only on
$h_K^2(-J_0\dot{z})$ integrated over time, which depends on the
multiset of velocities and their durations, not their ordering.

Iterating for all repeated Reeb vectors yields $z_3$ with:
\begin{itemize}
\item $A(z_3) \geq A(z_2)$ (each merge preserves or increases action),
\item $I_K(z_3) = I_K(z_2) = T$ (rearrangement preserves $I_K$),
\item each Reeb vector appears in at most one contiguous time interval.
\end{itemize}

\paragraph{Step 4: Rescale to restore normalization.}
After Steps 2--3, the action may have changed slightly:
$A(z_3) \geq A(z_2) \geq A(z_1) \geq T - \varepsilon$.
We restore the normalization $A(z) = T$ by scaling all time intervals
proportionally.

Define the scaling factor
\[
  \beta = \frac{T}{A(z_3)} \in (0, 1 + O(\varepsilon/T)].
\]
Construct $z_4$ by \emph{multiplying each time interval by $\beta$}:
\begin{itemize}
\item Original: velocity $R_i$ on interval of length $\Delta t_i$
\item Scaled: velocity $R_i$ on interval of length $\beta \cdot \Delta t_i$
\end{itemize}
The velocities remain unchanged (still pure Reeb vectors in direction and magnitude).

\textbf{Consequences:}
\begin{align*}
  T^{(4)} &= \beta \sum \Delta t_i = \beta \cdot T = \frac{T^2}{A(z_3)}, \\
  A(z_4) &= \beta^2 \cdot A(z_3)
    \quad\text{(action scales as [time]$^2$ when velocities constant)} \\
  &= \left(\frac{T}{A(z_3)}\right)^2 \cdot A(z_3) = \frac{T^2}{A(z_3)} = T^{(4)}, \\
  I_K(z_4) &= \beta \cdot I_K(z_3)
    \quad\text{($I_K = \frac{1}{4}\int h_K^2(-J_0\dot{z})\,dt$ is linear in time)} \\
  &= \frac{T}{A(z_3)} \cdot T = \frac{T^2}{A(z_3)}.
\end{align*}
Hence $A(z_4) = T^{(4)}$ (normalization restored) and
$I_K(z_4) = T^{(4)}$.

Since $z$ was a dual minimizer with $I_K(z) = T$, feasibility gives
$I_K(z_4) \geq T$, hence $T^2/A(z_3) \geq T$, so $A(z_3) \leq T$.
Combined with $A(z_3) \geq T - \varepsilon$, as $\varepsilon \to 0$
we have $A(z_3) \to T$, so $\beta \to 1$ and $T^{(4)} \to T$.

\paragraph{Step 5: Compactness and limit.}
A simple curve is encoded by finite data:
\begin{itemize}
\item A permutation $\sigma \in S_F$ (facet visit order),
\item Time durations $(t_1,\ldots,t_F) \in \mathbb{R}^F$ with $t_i \geq 0$,
      $\sum t_i = T^{(4)}$, $\sum t_i R_{\sigma(i)} = 0$ (closure).
\end{itemize}

% [AGENT-WRITTEN ARGUMENT, UNREVIEWED BY JÖRN]
% QC: uses finite-dimensional compactness (S_F × simplex) and the action
%   identity to bypass W^{1,2} convergence. A and I_K are computed
%   directly from finite data (σ, durations), so convergence of durations
%   gives convergence of A and I_K without functional analysis.
Each curve $z_{4,n}$ (from Step~4 with $\varepsilon_n \to 0$)
is determined by a permutation $\sigma^n \in S_F$
and durations $(t^n_1, \ldots, t^n_F) \in \mathbb{R}_{\geq 0}^F$
with $\sum t^n_i = T^{(4)}_n$ and $\sum t^n_i R_{\sigma^n(i)} = 0$.

The normalized durations
$(t^n_1/T^{(4)}_n, \ldots, t^n_F/T^{(4)}_n)$
lie in the $(F{-}1)$-simplex
$\Delta^{F-1} = \{s \in \mathbb{R}_{\geq 0}^F : \sum s_i = 1\}$.
Since $S_F \times \Delta^{F-1}$ is compact (finite $\times$ compact),
a subsequence converges:
\begin{itemize}
\item $\sigma^n = \sigma$ (constant, by pigeonhole),
\item $(t^n_1/T^{(4)}_n, \ldots, t^n_F/T^{(4)}_n)
      \to (s^\infty_1, \ldots, s^\infty_F) \in \Delta^{F-1}$,
\item $T^{(4)}_n \to T$ (from Step~4).
\end{itemize}
Define $z_*$ as the piecewise affine curve with ordering $\sigma$,
velocities $R_{\sigma(i)}$, and durations $t^\infty_i = s^\infty_i \cdot T$.
The closure condition
$\sum t^\infty_i R_{\sigma(i)}
  = \lim_n \sum t^n_i R_{\sigma(i)} = 0$
holds by continuity.
By the action identity~\eqref{eq:action-identity},
$A$ is a polynomial in the durations (for fixed $\sigma$), so
\[
  A(z_*) = \lim_{n\to\infty} A(z_{4,n})
  = \lim_{n\to\infty} T^{(4)}_n = T.
\]
Since all velocities are pure Reeb vectors,
$I_K(z_*) = T$ (same computation as Step~2).
As $z$ was a dual minimizer with $I_K(z) = T$,
and $z_*$ is feasible (closed, $A(z_*) = T$, centered
by translation) with $I_K(z_*) = T$,
the curve $z_*$ is also a dual minimizer.
It is simple by construction
(piecewise affine, pure Reeb velocities, each facet used at most once).

\paragraph{Wrapup.}
By Theorem~\ref{thm:primal-dual-equivalence}, the dual minimizer~$z_*$
corresponds to a primal minimizer $\gamma = z_* + b$ for some constant
$b \in \mathbb{R}^4$. Since $z_*$ is simple (piecewise linear with pure
Reeb vectors, each used at most once), $\gamma$ is also simple
(Definition~\ref{def:simple-reeb-orbit}). This completes the proof.
\end{proof}

